\clearpage\pagestyle{otMain}\setlength{\otBlocEspacement}{2.5mm}\fontsize{11}{15}\selectfont\corrigesChapitre{4}{Diviser : partages et restes}
\begin{blocLivre}{cor-4-1}\begin{corrige}{1}{Les premiers multiples}\demarche{Compter de 8 en 8.}
On compte de 8 en 8 : $\res{8,\ 16,\ 24,\ 32,\ 40}$. Chaque multiple s’écrit $8\times$ un entier.
\end{corrige}\end{blocLivre}
\begin{blocLivre}{cor-4-2}\begin{corrige}{2}{Pair ou impair}\demarche{Observer uniquement le chiffre des unités.}
On regarde le chiffre des unités. Pairs : $\res{346\text{ et }90}$. Impairs : $\res{1\,275\text{ et }83}$.
\end{corrige}\end{blocLivre}
\begin{blocLivre}{cor-4-3}\begin{corrige}{3}{Multiples de 5 et de 10}\demarche{Appliquer les critères des unités.}
Terminés par 0 ou 5 : $\res{85,\ 120,\ 905}$ sont multiples de 5. Seul $\res{120}$ finit par 0 : c’est le seul multiple de 10.
\end{corrige}\end{blocLivre}
\begin{blocLivre}{cor-4-4}\begin{corrige}{4}{Une division posée}\demarche{Encadrer 87 entre deux multiples de 4.}
Le quotient est $\res{21}$ et le reste $\res3$ : $\res{87=4\times21+3}$, avec $3<4$.
\end{corrige}\end{blocLivre}
\begin{blocLivre}{cor-4-5}\begin{corrige}{5}{Diviseur ou pas ?}\demarche{Chercher un produit exact.}
$7\times13=91$ : le reste est nul, donc $\res{7\text{ divise }91}$. $75=6\times12+3$ : le reste n’est pas nul, donc $\res{6\text{ ne divise pas }75}$.
\end{corrige}\end{blocLivre}
\begin{blocLivre}{cor-4-6}\begin{corrige}{6}{Partager équitablement}\demarche{Reconnaître la recherche de la valeur d’une part.}
$345:5=\res{69}$ billes chacun. Contrôle : $69\times5=345$, le partage tombe juste : 5 est un diviseur de 345.
\end{corrige}\end{blocLivre}
\begin{blocLivre}{cor-4-7}\begin{corrige}{7}{Quotient et contrôle}\demarche{Ne pas oublier le zéro intercalé au quotient.}
Le quotient est $\res{209}$ et le reste $\res4$. Contrôle : $6\times209+4=1\,258$ et $4<6$. Le zéro du quotient vient du chiffre 5, trop petit pour 6 après le premier partage.
\end{corrige}\end{blocLivre}
\begin{blocLivre}{cor-4-8}\begin{corrige}{8}{Encadrer par des multiples}\demarche{Chercher dans la table de 7.}
$7\times7=49$ et $7\times8=56$ : $\res{49<50<56}$. Cet encadrement donne le quotient de la division de 50 par 7 : c’est 7.
\end{corrige}\end{blocLivre}
\begin{blocLivre}{cor-4-9}\begin{corrige}{9}{Un reste trop grand}\demarche{Vérifier que le reste est plus petit que le diviseur.}
Son égalité est exacte, mais 10 dépasse le diviseur 8 : ce n’est pas une division euclidienne, on peut encore faire un paquet de 8. La bonne écriture est $\res{74=8\times9+2}$, avec $2<8$.
\end{corrige}\end{blocLivre}
\begin{blocLivre}{cor-4-10}\begin{corrige}{10}{Moitiés et quarts}\demarche{Diviser par 2, une ou deux fois.}
La moitié : $470:2=\res{235}$. Le quart, c’est la moitié de la moitié : $232:2=116$ puis $116:2=\res{58}$.
\end{corrige}\end{blocLivre}
\begin{blocLivre}{cor-4-11}\begin{corrige}{11}{Les œufs en boîtes}\demarche{Interpréter le reste selon la question posée.}
$230=6\times38+2$. On remplit $\res{38}$ boîtes complètes et il reste 2 œufs. Pour tout emporter, une boîte de plus est nécessaire : $\res{39}$ boîtes. La question choisit entre le quotient et le quotient augmenté de 1.
\end{corrige}\end{blocLivre}
\begin{blocLivre}{cor-4-12}\begin{corrige}{12}{Le vide-grenier}\demarche{Effectuer la division et observer le reste.}
$91:7=\res{13}$ livres par lot. Le reste est nul : $\res{\text{le partage tombe juste}}$, 7 est un diviseur de 91.
\end{corrige}\end{blocLivre}
\begin{blocLivre}{cor-4-13}\begin{corrige}{13}{Des équipes de huit}\demarche{Lire le quotient, puis raisonner sur le reste.}
$134=8\times16+6$ : on forme $\res{16}$ équipes complètes et $\res6$ élèves restent. Il manque $8-6=\res2$ élèves pour compléter une équipe de plus. Contrôle : $8\times17=136=134+2$.
\end{corrige}\end{blocLivre}
\begin{blocLivre}{cor-4-14}\begin{corrige}{14}{Le nombre mystère}\demarche{Lister les candidats puis appliquer chaque condition.}
Les multiples de 5 entre 82 et 94 sont 85 et 90. Seul 90 est pair : la réponse est $\res{90}$. Les deux conditions éliminent tous les autres nombres de l’intervalle.
\end{corrige}\end{blocLivre}
\begin{blocLivre}{cor-4-15}\begin{corrige}{15}{Toujours vrai ?}\demarche{Dépasser les exemples par un argument sur les paquets.}
$26+58=\res{84}$ et $134+72=\res{206}$ : deux sommes paires. Chaque nombre pair est fait de paquets de deux ; en réunissant les paquets, la somme reste faite de paquets de deux : l’affirmation est $\res{\text{toujours vraie}}$.
\end{corrige}\end{blocLivre}
